\section{Valuative covering by finite trees of blowup charts}

\begin{definition}[Finite trees of affine charts]\label{def:sc0-rees-chart-tree}
\lean{EnvelopeProperty.reesTreeChart,EnvelopeProperty.ReesChartTree,EnvelopeProperty.ReesChartTree.leaves}
\leanok
Let $B$ be a commutative ring, $K$ a field with a $B$-algebra structure, and $C\subseteq K$ a
$B$-subalgebra. For a finite family $g_1,\ldots,g_n\in C$ and $g_i\ne0$, put
\[
 C[g/g_i]=C[g_1/g_i,\ldots,g_n/g_i]\subseteq K.
\]
A finite chart tree rooted at $C$ is either the leaf $C$, or a nonzero finite family $g$ at $C$
together with a finite chart tree rooted at $C[g/g_i]$ for each $g_i\ne0$.
The leaves of a node are the union of the leaves of its children.
\end{definition}

\begin{lemma}[Principalization on an affine chart]\label{lem:sc0-rees-chart-principal}
\uses{def:sc0-rees-chart-tree}
\lean{EnvelopeProperty.reesTreeChartIncl_mul_ratio,EnvelopeProperty.reesTreeChart_map_span_eq}
\leanok
If $g_i\ne0$, the extension of the ideal $(g_1,\ldots,g_n)\subseteq C$ to $C[g/g_i]$ equals
the principal ideal generated by $g_i$.
\end{lemma}
\begin{proof}
Every $g_j$ equals $g_i(g_j/g_i)$ in the chart algebra, so the extended ideal is contained
in $(g_i)$. Conversely, $g_i$ is one of the original generators and therefore belongs to the
extended ideal. These two inclusions prove the equality.
\end{proof}

\begin{lemma}[A chart contained in a valuation ring]\label{lem:sc0-rees-chart-in-valuation-ring}
\uses{def:sc0-rees-chart-tree}
\lean{EnvelopeProperty.exists_reesTreeChart_le}
\leanok
Let $V\subseteq K$ be a valuation ring containing $C$, and let $g$ be a nonzero finite family in
$C$. There is an index $i$ with $g_i\ne0$ and $C[g/g_i]\subseteq V$.
\end{lemma}
\begin{proof}
Let $v$ be the multiplicative valuation associated with $V$, so $x\in V$ if and only if $v(x)\le1$.
Choose $i$ maximizing $v(g_i)$. Some $g_j$ is nonzero, hence $v(g_i)\ge v(g_j)>0$, so $g_i\ne0$.
For every $j$, $v(g_j)\le v(g_i)$, which implies $g_j/g_i\in V$. Since $V$ contains $C$ and is
closed under ring operations, it contains the subalgebra generated by these ratios.
\end{proof}

\begin{lemma}[The leaves form a finite valuative cover]\label{lem:sc0-rees-tree-cover}
\uses{def:sc0-rees-chart-tree,lem:sc0-rees-chart-in-valuation-ring}
\lean{EnvelopeProperty.ReesChartTree.finite_leaves,EnvelopeProperty.ReesChartTree.leaves_nonempty,EnvelopeProperty.ReesChartTree.exists_leaf_le}
\leanok
Every finite chart tree has a finite, nonempty set of leaves. If a valuation ring $V\subseteq K$
contains its root, it contains one of its leaves.
\end{lemma}
\begin{proof}
Proceed by induction on the tree. For a leaf all three assertions hold immediately.
For a node there are finitely many children, so their finite leaf sets have finite union.
The defining nonzero-family condition gives a child, whose leaf set is nonempty by induction.
If $V$ contains the root, Lemma~\ref{lem:sc0-rees-chart-in-valuation-ring} gives a child whose root
is contained in $V$. Apply induction to that child.
\end{proof}

\begin{lemma}[Containment and finite type of leaves]\label{lem:sc0-rees-tree-finite-type}
\uses{def:sc0-rees-chart-tree}
\lean{EnvelopeProperty.ReesChartTree.root_le_of_mem_leaves,EnvelopeProperty.ReesChartTree.fg_of_mem_leaves,EnvelopeProperty.ReesChartTree.finiteType_of_mem_leaves}
\leanok
Each leaf contains the root algebra. If the root is finitely generated over $B$, so is every leaf.
\end{lemma}
\begin{proof}
Each child algebra is generated by the parent algebra together with finitely many coordinate
ratios. Consequently it contains the parent, and the union of a finite set of parent generators
with these ratios generates the child. Induction on the tree proves both assertions.
\end{proof}
